\chapter{The Dirac Gate in the Compiler}
\label{chap:dirac-gate-compiler}

The experimental construction writes the Dirac equation as physics. The compiler reads it as a gate: a
typed passage through which a state may move only when the local algebra, the transport rule, and the mass
term all agree. That distinction matters. The compiler does not need an interpretation of spinor motion in
order to check a Dirac step; it needs the shape of the step as a judgment. The spinor is a value in a
module, the connection is the environment against which neighboring values are compared, and the gamma
matrices are not decoration but the algebraic certificate that the first-order operator squares against the
metric it was promised.

A first-order equation is exactly the kind of thing a compiler can see. It does not ask the machine to
look across the whole path at once. It asks for one admissible comparison at one address, then another, then
another, each comparison carrying enough structure for the next one to be checked. Written in the familiar
compressed form,
\[
  D_A\psi
    =
  \gamma^\mu(\nabla_\mu^A \psi) + m\psi ,
\]
the expression is already a compiler contract. The symbol \(A\) tells the checker what background transport
is in force; \(\nabla_\mu^A\) tells it how a spinor may be compared after displacement; the Clifford law
\[
  \gamma^\mu\gamma^\nu+\gamma^\nu\gamma^\mu
    =
  2g^{\mu\nu}I
\]
tells it why the first-order gate is compatible with the second-order geometry beneath it. The equation is
not only an equality of fields. It is a discipline for moving a value through context without losing the
metric that made the context meaningful.

\section{First order means locally checkable}
\label{se:first-order-locally-checkable}

The compiler's affection for first-order structure is not aesthetic. A first-order operator exposes its
obligations at the same scale as elaboration. A term enters with a type, an environment, and a local
neighborhood of dependencies. If the operator demanded a global minimizer at every use, the checker would
have to solve an entire variational problem before admitting a single step. The Dirac gate does something
more useful: it records the differential comparison, the algebra that controls that comparison, and the
residual left after the comparison is made.

In that sense the Dirac equation is not remote from the earlier laws of the compiler. Poison was a law of
what cannot pass through an altered context. The Dirac gate is a positive version of the same phenomenon:
it says exactly what extra structure must be present for passage through a changing context to remain
well-typed. A bare derivative compares values that live in different fibers as though they were already in
one place. A covariant derivative refuses that shortcut. It brings the comparison map with it. The compiler
recognizes the refusal: no value may be compared across a boundary unless the boundary also supplies the
transport that makes the comparison legal.

\section{Gamma as a typing law}
\label{se:gamma-as-typing-law}

The gamma matrices are often introduced as analytic machinery, but the compiler reads them as a typing law
for directions. Directional derivatives by themselves form a list of partial motions. Clifford
multiplication binds that list into a single operator whose square remembers the metric. Without the
Clifford relation, the first-order expression would still have syntax; it would not have the promised
geometry. The checker therefore treats the gamma algebra the way it treats a class law or a record field:
not as a computation to be rediscovered at each use, but as a certificate carried by the structure that
licenses the computation.

This is why the Dirac operator belongs naturally beside the elaborator rather than outside it. The
elaborator already knows the difference between a term that merely parses and a term whose implicit
structure resolves. A spinor expression with an unlicensed gamma action is only syntax. A spinor expression
whose Clifford action is carried in context can be elaborated into a judgment. The mathematical content has
not been reduced to bookkeeping; the bookkeeping is the visible form of the content at the scale where the
compiler can enforce it.

\section{A gate, not an oracle}
\label{se:dirac-gate-not-oracle}

The Dirac gate gives the compiler a way to admit or reject local motion, not a way to conjure a global
solution. A global solution still requires a variational reading, a topology for admissible states, and a
norm in which residuals can be measured. That is where the weak equation enters. The point of starting with
the gate is that the weak equation should not appear as a retreat from exactness. It is a change in where
exactness is demanded. The local algebra remains exact. The transport rule remains exact. The residual is
not ignored; it is moved into a space where the compiler can ask the only question it can answer honestly:
which finite family of tests has been paid for, and what residual do those tests read?

Read this way, the Dirac gate is the compiler-facing boundary of the analytic apparatus. On one side are
spinors, transport, gamma laws, mass, and charge. On the other side are judgments, contexts, obligations,
and residuals. The gate matters because it preserves both languages at once. It does not explain the
physics away; it gives the compiler a typed opening through which the physics can be measured.
